% ===============================================================
% Document Class ================================================
\documentclass{article}

% ===============================================================
% Graphic Packages ==============================================
\usepackage{graphicx}       % Permite a inclusão de imagens no documento.
\usepackage{tocloft} % Pacote para customização de listas de conteúdos, figuras e tabelas
\usepackage{tikz}           % Ferramenta poderosa para criar gráficos programaticamente dentro do LaTeX.
\usetikzlibrary{calc}       % Extensão da biblioteca TikZ que permite cálculos mais complexos de coordenadas.
\usepackage[portuguese]{babel}
\usepackage{tocloft} % Pacote para customização de listas
\renewcommand{\listfigurename}{Lista de Figuras} % Altera o título para português
\AtBeginDocument{\renewcommand{\contentsname}{Sumário}} % Altera nome de conteudo para sumário
% ===============================================================
% Mathematical Tools ============================================
\usepackage{amsmath}        % Melhora a aparência e a flexibilidade de comandos matemáticos.
\usepackage{siunitx}        % Facilita o uso de unidades do Sistema Internacional e ajuda a formatar números complexos.

% ===============================================================
% Font and Text Appearance ======================================
\usepackage{mathptmx}       % Altera a fonte padrão do documento para Times New Roman.

% ===============================================================
% Table of Contents Customization ===============================
\usepackage{tocloft}  % Oferece controle total sobre a aparência das listas de conteúdos, figuras, tabelas, etc.

% ===============================================================
% Figure Positioning ============================================
\usepackage{float}     % Melhora a interface para definir o posicionamento de objetos flutuantes como figuras e tabelas.

% ===============================================================
% Paragraph Spacing and Indentation =============================
\usepackage{setspace}  % Permite o ajuste fino do espaçamento entre linhas.
\usepackage{indentfirst} % Adiciona indentação ao primeiro parágrafo de cada seção.

% ===============================================================
% Page Layout ===================================================
\usepackage[a4paper, top=3cm, bottom=2cm, left=3cm, right=2cm]{geometry}  % Define as margens de todo o documento.
\setlength{\parindent}{4em}  % Define o tamanho da indentação para todos os parágrafos.
\setlength{\emergencystretch}{3em}

% ===============================================================
% Section Heading Customization =================================
\makeatletter
\renewcommand\paragraph{\@startsection{paragraph}{4}{\z@}%
    {2ex plus 1ex minus .2ex}%
    {1em}%
    {\normalfont\normalsize\bfseries}}
\makeatother

% ===============================================================
% References =================================
\usepackage{cite} 


% ===============================================================
% Urls ==========================================================
\usepackage{hyperref} % Para criar links clicáveis
\hypersetup{
    colorlinks=true,
    linkcolor=blue,
    filecolor=magenta,
    urlcolor=blue,
    citecolor=blue,
    pdfborder={0 0 0},  % Remove o quadrado ao redor dos links
    breaklinks=true  % Permite que os links sejam quebrados entre as linhas
}
\usepackage{xurl} % Melhora a quebra de links em qualquer lugar

% ===============================================================
% Bloco de código ===============================================
\usepackage{listings}
\usepackage{xcolor}
\usepackage{courier}
\lstset{
  backgroundcolor=\color[RGB]{249,246,239},
  basicstyle=\ttfamily\footnotesize,
  breaklines=true,
  frame=single,
  numbers=left,
  numberstyle=\tiny\color{gray},
  keywordstyle=\color[RGB]{40,40,255},
  commentstyle=\color[RGB]{0,125,0},
  stringstyle=\color[RGB]{255,0,0},
  showstringspaces=false,
  rulecolor=\color{black},
  captionpos=b,
  abovecaptionskip=5pt,
  belowcaptionskip=5pt,
  xleftmargin=0.15\textwidth,
  xrightmargin=0.15\textwidth,
  morecomment=[l]{//}
}
\renewcommand{\lstlistlistingname}{Lista de Códigos} % Comando para renomear o título da Lista de Códigos

% ===============================================================
% Begin Document ================================================
\begin{document}

% Capa basica do Ratamero
\begin{flushright}
  \textbf{Trabalho da Disciplina Estudo de Casos Empresariais} \\
  Orientador: Prof. Leandro de Amorim Ratamero
\end{flushright}

% Título da atividade
{\large\bfseries Trabalho 7 - Ifood \par}

% Nome do estudante
{\large Guilherme Cagide Fialho \par}

% ===============================================================
% Capa ==========================================================
% \begin{titlepage}
    \centering
    % Desenhar a margem
    \begin{tikzpicture}[remember picture, overlay]
        \draw[line width = 4pt] ($(current page.north west) + (20mm, -20mm)$) rectangle ($(current page.south east) + (-10mm,10mm)$);
        \draw[line width = 1pt] ($(current page.north west) + (21.5mm,-21.5mm)$) rectangle ($(current page.south east) + (-11.5mm,11.5mm)$);
    \end{tikzpicture}

    % Linha de logos
    \includegraphics[height=0.151\textwidth]{header/contra-capa/assets/uerj.png}\hfill
    \includegraphics[height=0.15\textwidth]{header/contra-capa/assets/iprj.jpeg}\hfill

    \vspace{2cm} % Espaço vertical após os logos

    % Informações do curso
    {\Large\bfseries Instituto Politécnico do Estado do Rio de Janeiro \par}
    \vspace{0.5cm}
    {\large Curso de Engenharia da Computação \par}

    \vspace{3cm} % Espaço vertical antes do título da atividade

    % Nome do estudante
    {\large\bfseries Guilherme Cagide Fialho \par}

    \vspace{3cm}

    % Título da atividade
    {\large\bfseries Trabalho 2 - Microsoft\par}

    \vfill % Empurra o restante para o fundo da página

    % Local e ano
    {\large\bfseries Nova Friburgo \par}
    \vspace{0.3cm}
    {\large\bfseries 2024 \par}
\end{titlepage}
\newpage % Começa uma nova página após a capa


% ===============================================================
% Contra Capa ===================================================
% \begin{titlepage}
    \centering
    % Desenhar a margem
    \begin{tikzpicture}[remember picture, overlay]
        \draw[line width = 4pt] ($(current page.north west) + (20mm, -20mm)$) rectangle ($(current page.south east) + (-10mm,10mm)$);
        \draw[line width = 1pt] ($(current page.north west) + (21.5mm,-21.5mm)$) rectangle ($(current page.south east) + (-11.5mm,11.5mm)$);
    \end{tikzpicture}

    % Linha de logos
    \includegraphics[height=0.151\textwidth]{header/contra-capa/assets/uerj.png}\hfill
    \includegraphics[height=0.15\textwidth]{header/contra-capa/assets/iprj.jpeg}\hfill

    \vspace{2cm} % Espaço vertical após os logos

    % Informações do curso
    {\Large\bfseries Instituto Politécnico do Estado do Rio de Janeiro \par}
    \vspace{0.5cm}
    {\large Graduação em Engenharia da Computação \par}

    \vspace{3cm} % Espaço vertical antes do título da atividade

    % Nome do estudante
    {\large Guilherme Cagide Fialho \par}

    \vspace{1.5cm}

    % Título da atividade
    {\large\bfseries Trabalho 5 - Amazon \par}

    \vspace{1cm} % Espaço vertical antes do título da atividade

    % Informações do projeto
    \begin{flushright}
        \begin{minipage}{0.5\textwidth}
            \large
            \raggedleft % Garante que o texto dentro da minipage seja alinhado à direita
            Trabalho da Disciplina Estudo de Casos Empresariais
        \end{minipage}
    \end{flushright}

    \vspace{1.5cm}

    % Informações do orientador
    \begin{flushleft}
        \begin{minipage}{0.5\textwidth}
            \large
            \raggedright
            Orientador: Prof. Leandro de Amorim Ratamero
        \end{minipage}
    \end{flushleft}


    \vfill % Empurra o restante para o fundo da página

    % Local e ano
    {\large Nova Friburgo \par}
    \vspace{0.3cm}
    {\large 2024 \par}
\end{titlepage}
\newpage % Começa uma nova página após a capa


% ===============================================================
% Resumo ========================================================
% \begin{titlepage}
    \thispagestyle{empty} % Remove números de página
    \setstretch{1.5} % Espaçamento entre linhas, certifique-se de que o pacote setspace está incluído em document.tex

    \begin{center}
        \textbf{\Large RESUMO}
    \end{center}


    \vspace{1cm} % Espaço vertical

    \noindent CAGIDE FIALHO, G. Relatório do projeto de Modelagem
    e controle de sistemas. 2024. 54 f. Trabalho de Conclusão de Disciplina Modelagem e Controle de Sistemas (Graduação em
    Engenharia da computação) – Graduação em Engenharia da Computação, Universidade
    do Estado do Rio de Janeiro, Nova Friburgo, 2024.

    \vspace{0.4cm} % Espaço vertical

    Este trabalho explora a modelagem e o controle de um sistema dinâmico do tipo massa-mola-amortecedor, utilizando a plataforma Scilab para desenvolvimento e simulação. O foco do estudo está na implementação de modelos matemáticos para descrever a dinâmica do sistema e na análise de sua resposta sob diversas condições iniciais, sem a aplicação de forças externas. Utilizando a ferramenta Xcos, um componente gráfico do Scilab, realizamos simulações que permitiram uma análise visual e quantitativa das respostas transientes do sistema. O estudo destaca a influência dos parâmetros físicos, como a massa, o coeficiente de amortecimento, e a constante da mola, nas características de resposta do sistema. Além disso, técnicas de controle foram empregadas para ajustar a resposta do sistema, demonstrando como o amortecimento pode contribuir para a estabilização após perturbações e enfatizando a relevância de uma parametrização cuidadosa para alcançar um comportamento eficaz do sistema. Este projeto contribui para a compreensão das teorias de controle aplicáveis em sistemas mecânicos e outros contextos de sistemas dinâmicos na engenharia.
    \vspace{0.4cm} % Espaço vertical

    \textbf{Palavras-chave}: Modelagem e Controle, Sistema Massa-Mola-Amortecedor, Simulação, Scilab, Xcos.
\end{titlepage}


% ===============================================================
% Abstract ======================================================
% \begin{titlepage}
    \thispagestyle{empty} % Remove page numbers
    \setstretch{1.5} % Line spacing, make sure the setspace package is included in document.tex

    \begin{center}
        \textbf{\Large ABSTRACT}
    \end{center}

    \vspace{1cm} % Vertical space

    \noindent CAGIDE FIALHO, G. Report on the Modeling and Control Systems Project. 2024. 54 p. Final Project for the Modeling and Control Systems Course (Bachelor’s degree in Computer Engineering) - Bachelor’s Degree in Computer Engineering, State University of Rio de Janeiro, Nova Friburgo, 2024.

    \vspace{0.4cm} % Vertical space

    This work explores the modeling and control of a dynamic system of the mass-spring-damper type, using the Scilab platform for development and simulation. The focus of the study is on implementing mathematical models to describe the system dynamics and analyzing its response under various initial conditions without the application of external forces. Using Xcos, a graphical component of Scilab, we conducted simulations that allowed for visual and quantitative analysis of the system's transient responses. The study highlights the influence of physical parameters, such as mass, damping coefficient, and spring constant, on the system's response characteristics. Furthermore, control techniques were employed to adjust the system's response, demonstrating how damping can contribute to stabilization after disturbances and emphasizing the importance of careful parameterization to achieve effective system behavior. This project contributes to the understanding of control theories applicable to mechanical systems and other dynamic system contexts in engineering.
    \vspace{0.4cm} % Vertical space

    \textbf{Keywords}: Modeling and Control, Mass-Spring-Damper System, Simulation, Scilab, Xcos.
\end{titlepage}


% ===============================================================
% Lista de figuras ==============================================
% \newpage
\setstretch{1.5} % Espaçamento entre linhas, certifique-se de que o pacote setspace está incluído em document.tex

% % Estilo de citação (plain, abbrv, ieeetr, etc.)
% \bibliographystyle{IEEEtran}  % Estilo que inclui URLs
% % \bibliographystyle{plain}

% \bibliography{
%     perguntas/1-foundation/latex/references,
%     perguntas/2-graduation/latex/references,
%     perguntas/3-main-productions/latex/references,
%     perguntas/4-location/latex/references,
%     perguntas/5-competitors/latex/references,
%     perguntas/6-protection/latex/references
% }

% Estilo de citação (plain, abbrv, ieeetr, etc.)
\bibliographystyle{IEEEtran}  % Estilo que inclui URLs
% \bibliographystyle{plain}

\bibliography{perguntas/1-foundation/latex/references,perguntas/2-graduation/latex/references,perguntas/3-main-products/latex/references.bib,perguntas/4-location/latex/references,perguntas/5-competitors/latex/references,perguntas/6-protection/latex/references, perguntas/7-profit/latex/references, perguntas/8-train/latex/references, perguntas/9-social-programs/latex/references, perguntas/10-green-purposes/latex/references}


% ===============================================================
% Lista de códigos ==============================================
% \newpage
\setstretch{1.5} % Espaçamento entre linhas, certifique-se de que o pacote setspace está incluído em document.tex

% % Estilo de citação (plain, abbrv, ieeetr, etc.)
% \bibliographystyle{IEEEtran}  % Estilo que inclui URLs
% % \bibliographystyle{plain}

% \bibliography{
%     perguntas/1-foundation/latex/references,
%     perguntas/2-graduation/latex/references,
%     perguntas/3-main-productions/latex/references,
%     perguntas/4-location/latex/references,
%     perguntas/5-competitors/latex/references,
%     perguntas/6-protection/latex/references
% }

% Estilo de citação (plain, abbrv, ieeetr, etc.)
\bibliographystyle{IEEEtran}  % Estilo que inclui URLs
% \bibliographystyle{plain}

\bibliography{perguntas/1-foundation/latex/references,perguntas/2-graduation/latex/references,perguntas/3-main-products/latex/references.bib,perguntas/4-location/latex/references,perguntas/5-competitors/latex/references,perguntas/6-protection/latex/references, perguntas/7-profit/latex/references, perguntas/8-train/latex/references, perguntas/9-social-programs/latex/references, perguntas/10-green-purposes/latex/references}


% ===============================================================
% Lista de siglas e abreviaturas ================================
% \newpage
\setstretch{1.5} % Espaçamento entre linhas, certifique-se de que o pacote setspace está incluído em document.tex

% % Estilo de citação (plain, abbrv, ieeetr, etc.)
% \bibliographystyle{IEEEtran}  % Estilo que inclui URLs
% % \bibliographystyle{plain}

% \bibliography{
%     perguntas/1-foundation/latex/references,
%     perguntas/2-graduation/latex/references,
%     perguntas/3-main-productions/latex/references,
%     perguntas/4-location/latex/references,
%     perguntas/5-competitors/latex/references,
%     perguntas/6-protection/latex/references
% }

% Estilo de citação (plain, abbrv, ieeetr, etc.)
\bibliographystyle{IEEEtran}  % Estilo que inclui URLs
% \bibliographystyle{plain}

\bibliography{perguntas/1-foundation/latex/references,perguntas/2-graduation/latex/references,perguntas/3-main-products/latex/references.bib,perguntas/4-location/latex/references,perguntas/5-competitors/latex/references,perguntas/6-protection/latex/references, perguntas/7-profit/latex/references, perguntas/8-train/latex/references, perguntas/9-social-programs/latex/references, perguntas/10-green-purposes/latex/references}


% ===============================================================
% Materiais utilizados ==========================================
% \newpage
\setstretch{1.5} % Espaçamento entre linhas, certifique-se de que o pacote setspace está incluído em document.tex

% % Estilo de citação (plain, abbrv, ieeetr, etc.)
% \bibliographystyle{IEEEtran}  % Estilo que inclui URLs
% % \bibliographystyle{plain}

% \bibliography{
%     perguntas/1-foundation/latex/references,
%     perguntas/2-graduation/latex/references,
%     perguntas/3-main-productions/latex/references,
%     perguntas/4-location/latex/references,
%     perguntas/5-competitors/latex/references,
%     perguntas/6-protection/latex/references
% }

% Estilo de citação (plain, abbrv, ieeetr, etc.)
\bibliographystyle{IEEEtran}  % Estilo que inclui URLs
% \bibliographystyle{plain}

\bibliography{perguntas/1-foundation/latex/references,perguntas/2-graduation/latex/references,perguntas/3-main-products/latex/references.bib,perguntas/4-location/latex/references,perguntas/5-competitors/latex/references,perguntas/6-protection/latex/references, perguntas/7-profit/latex/references, perguntas/8-train/latex/references, perguntas/9-social-programs/latex/references, perguntas/10-green-purposes/latex/references}


% ===============================================================
% Sumário =======================================================
% \newpage
\setstretch{1.5} % Espaçamento entre linhas, certifique-se de que o pacote setspace está incluído em document.tex

% % Estilo de citação (plain, abbrv, ieeetr, etc.)
% \bibliographystyle{IEEEtran}  % Estilo que inclui URLs
% % \bibliographystyle{plain}

% \bibliography{
%     perguntas/1-foundation/latex/references,
%     perguntas/2-graduation/latex/references,
%     perguntas/3-main-productions/latex/references,
%     perguntas/4-location/latex/references,
%     perguntas/5-competitors/latex/references,
%     perguntas/6-protection/latex/references
% }

% Estilo de citação (plain, abbrv, ieeetr, etc.)
\bibliographystyle{IEEEtran}  % Estilo que inclui URLs
% \bibliographystyle{plain}

\bibliography{perguntas/1-foundation/latex/references,perguntas/2-graduation/latex/references,perguntas/3-main-products/latex/references.bib,perguntas/4-location/latex/references,perguntas/5-competitors/latex/references,perguntas/6-protection/latex/references, perguntas/7-profit/latex/references, perguntas/8-train/latex/references, perguntas/9-social-programs/latex/references, perguntas/10-green-purposes/latex/references}


% ===============================================================
% Introdução ====================================================
% \newpage
\setstretch{1.5} % Espaçamento entre linhas, certifique-se de que o pacote setspace está incluído em document.tex

% % Estilo de citação (plain, abbrv, ieeetr, etc.)
% \bibliographystyle{IEEEtran}  % Estilo que inclui URLs
% % \bibliographystyle{plain}

% \bibliography{
%     perguntas/1-foundation/latex/references,
%     perguntas/2-graduation/latex/references,
%     perguntas/3-main-productions/latex/references,
%     perguntas/4-location/latex/references,
%     perguntas/5-competitors/latex/references,
%     perguntas/6-protection/latex/references
% }

% Estilo de citação (plain, abbrv, ieeetr, etc.)
\bibliographystyle{IEEEtran}  % Estilo que inclui URLs
% \bibliographystyle{plain}

\bibliography{perguntas/1-foundation/latex/references,perguntas/2-graduation/latex/references,perguntas/3-main-products/latex/references.bib,perguntas/4-location/latex/references,perguntas/5-competitors/latex/references,perguntas/6-protection/latex/references, perguntas/7-profit/latex/references, perguntas/8-train/latex/references, perguntas/9-social-programs/latex/references, perguntas/10-green-purposes/latex/references}


% ===============================================================
% Atividades ====================================================
\section{Qual é o nome da empresa em estudo? Quando foi fundada? Qual o(s) nome(s) do(s) fundador(es)?}

A empresa que estamos analisando é a \textbf{Amazon Inc.}, uma gigante do setor de tecnologia que atua principalmente nas áreas de comércio eletrônico, computação em nuvem, streaming digital e inteligência artificial. A Amazon foi criada em \textbf{5 de julho de 1994}, na cidade de Bellevue, no estado de Washington, EUA. O responsável por sua fundação é \textbf{Jeffrey Preston Bezos}, mais conhecido como \textit{Jeff Bezos}, que naquela época tomou a corajosa decisão de deixar uma carreira estável e de sucesso em uma firma de investimentos para apostar em uma ideia revolucionária. O que começou como uma simples livraria online se transformou, ao longo dos anos, na maior varejista da internet e em uma das empresas mais valiosas e influentes do mundo, vendendo de tudo, de livros a serviços digitais avançados \cite{bezosbiography2024}. % Completa e Validada
\section{Qual é a formação de cada fundador?}

Os fundadores da iFood possuem formações diversas que contribuíram para o sucesso da empresa no setor de tecnologia e delivery de alimentos:

\begin{itemize}
    \item \textbf{Patrick Sigrist}: Engenheiro Florestal, com Mestrado em Análise de Dados Espaciais, além de possuir um MBA. Sigrist é um empreendedor em série, tendo fundado outras startups, como a Nomad e a Yellow Ventures \cite{sigrist2023}.
    
    \item \textbf{Felipe Fioravante}: Economista, com experiência prévia em consultoria. Antes de ingressar na iFood, ele deixou um emprego em uma empresa de consultoria para se aventurar no mundo das startups \cite{fioravante2023}.
    
    \item \textbf{Eduardo Baer}: Possui formação em Administração de Empresas. Além de cofundador da iFood, ele também é envolvido em outros projetos no setor de tecnologia e investimentos \cite{baer2023}.
    
    \item \textbf{Guilherme Bonifácio}: Formado em Engenharia Civil, também com MBA em Gestão de Negócios. Bonifácio tem vasta experiência no setor de tecnologia e empreendedorismo, contribuindo para a estruturação da iFood desde seu início \cite{bonifacio2023}.
\end{itemize}

Esses fundadores trouxeram uma diversidade de formações e experiências, o que permitiu à iFood desenvolver uma plataforma robusta e inovadora no mercado de delivery.

 % Completa e Validadada
\section{Quais são os produtos e/ou serviços mais conhecidos desta empresa no mercado?}

A Petrobras é amplamente reconhecida pela sua expertise em exploração de petróleo em águas profundas e ultraprofundas, abrangendo todas as fases desde o estudo geológico até a exploração efetiva de petróleo e gás. Essa exploração resulta em uma variedade de produtos derivados do petróleo:

\begin{itemize}
    \item \textbf{Gasolina:} Amplamente utilizada em veículos, é um dos principais produtos derivados do petróleo da Petrobras.
    \item \textbf{Óleo Diesel:} Essencial para o transporte e a indústria, a Petrobras oferece variantes que incluem misturas com conteúdo renovável.
    \item \textbf{Gás Liquefeito de Petróleo (GLP):} Conhecido como gás de cozinha, é usado extensivamente tanto em residências quanto na indústria.
    \item \textbf{Parafina:} Aplicada em diversos produtos, desde velas até cosméticos e borrachas.
    \item \textbf{Produtos Asfálticos:} Utilizados em projetos de pavimentação e infraestrutura civil.
\end{itemize}

Além desses produtos, a Petrobras é líder no desenvolvimento de tecnologias avançadas para exploração petrolífera em condições extremas, como as encontradas em reservatórios do pré-sal, e na produção de eteno e propeno, fundamentais para a indústria petroquímica. % Completa e validada
\section{Quais os países nos quais a empresa atua?}

Atualmente, o iFood atua apenas no Brasil, onde tem sua sede principal em São Paulo. A empresa opera em todos os estados brasileiros e é líder no setor de delivery de comida e produtos de mercado no país.

Anteriormente, a iFood teve operações internacionais em países como Colômbia, onde encerrou suas atividades em 2022, e México, onde atuou até 2020. Também já esteve presente na Argentina, mas vendeu sua operação local em 2018 \cite{ifood2023history}.

 % Completa e validada
\section{Quais são os principais concorrentes destas empresas no que se refere ao item 3?}

No mercado de streaming, que é o principal serviço oferecido pela \textbf{Netflix}, a empresa enfrenta uma forte concorrência global. Os principais concorrentes variam de acordo com a região e a oferta de conteúdo, mas incluem grandes empresas de tecnologia e comunicação. Entre os principais concorrentes estão:

\begin{itemize}
    \item \textbf{Amazon Prime Video}: Um dos principais concorrentes globais da Netflix, o Amazon Prime Video faz parte do ecossistema da Amazon e oferece uma vasta biblioteca de filmes, séries e conteúdos originais. Sua estratégia de integração com outros serviços do Amazon Prime fortalece sua competitividade.

    \item \textbf{Disney+}: O serviço de streaming da Disney rapidamente se tornou um dos maiores concorrentes da Netflix, graças ao enorme catálogo de produções da Disney, Pixar, Marvel, Star Wars e National Geographic. Além disso, a plataforma investe em produções originais e séries populares que atraem um público global.

    \item \textbf{HBO Max}: A HBO Max, parte do conglomerado WarnerMedia, é outro competidor importante. Seu catálogo inclui séries de prestígio, como \textit{Game of Thrones}, \textit{Succession} e produções de longa data da HBO, além de filmes de grandes estúdios e conteúdo exclusivo.

    \item \textbf{Apple TV+}: Embora tenha entrado mais recentemente no mercado de streaming, a Apple TV+ já se destaca por suas produções originais de alta qualidade, como \textit{Ted Lasso} e \textit{The Morning Show}. A integração com o ecossistema Apple torna o serviço um concorrente relevante, especialmente nos EUA.

    \item \textbf{Hulu}: Disponível principalmente nos Estados Unidos, o Hulu é outro forte concorrente, com um modelo de negócios que combina conteúdo sob demanda com transmissões ao vivo de canais de TV. Ele também oferece uma combinação de filmes, séries e conteúdo original.

    \item \textbf{Paramount+}: O serviço de streaming da Paramount oferece um catálogo vasto, incluindo séries populares, filmes e eventos esportivos ao vivo, além de ser integrado a outras redes como CBS.

\end{itemize}

Esses competidores buscam se diferenciar por meio de conteúdo exclusivo, integrações com outros serviços e preços competitivos. A competição no mercado de streaming é intensa e continuamente evolui à medida que novas plataformas e tecnologias surgem \cite{concorrentes_streaming}. % Completa e validada
\section{Há proteção intelectual do(s) produto(s) ou processo(s) mais famosos da empresa? Qual o tipo?}
A Amazon protege seus produtos e serviços mais conhecidos por meio de um vasto portfólio de patentes, marcas registradas e direitos autorais. Com mais de 33 mil patentes ativas globalmente, a empresa assegura inovações em áreas como computação em nuvem (AWS), inteligência artificial (Alexa), e processos de e-commerce, como o \textbf{Amazon Echo}, o \textbf{Kindle} e o \textbf{Prime}. Além disso, a Amazon detém mais de 1.700 marcas registradas mundialmente, cobrindo nomes como \textbf{Alexa}, \textbf{Kindle} e \textbf{AWS}, garantindo exclusividade no uso comercial. Direitos autorais protegem o conteúdo digital do \textbf{Kindle} e do \textbf{Prime Video}, assegurando controle sobre a distribuição de seus conteúdos literários e audiovisuais \cite{amazonPatents2024}. % Completa e validada
\section{Qual o lucro líquido anual mais recente dessa empresa?}

Em 2023, o Nubank alcançou um lucro líquido recorde de \textbf{US\$ 1 bilhão}, revertendo o prejuízo do ano anterior. Com uma receita anual superior a \textbf{US\$ 8 bilhões} e uma base de 93,9 milhões de clientes nos países onde opera (Brasil, México e Colômbia), o Nubank expandiu significativamente seus produtos e serviços financeiros, consolidando-se como uma das maiores fintechs da América Latina \cite{nubank2023resultados}.

 % Completo e Validado
\section{Esta empresa possui programa(s) de treinamento do seu pessoal? Quais?}

A Microsoft oferece alguns programas de treinamento para seus funcionários, visando o desenvolvimento de habilidades técnicas e a obtenção de certificações profissionais.

\begin{itemize}
    \item \textbf{Microsoft Learn}: Plataforma de aprendizado que disponibiliza cursos interativos, caminhos de aprendizagem e módulos que cobrem uma ampla gama de produtos e tecnologias Microsoft, incluindo áreas de nuvem e desenvolvimento de software.
    \item \textbf{Microsoft Virtual Training Days}: Eventos de capacitação técnica gratuitos que proporcionam aos participantes uma compreensão aprofundada das tecnologias da Microsoft Cloud, além de prepará-los para exames de certificação.
    \item \textbf{Certificações Microsoft}: Programas que validam as competências técnicas adquiridas em várias tecnologias da Microsoft, tais como Azure, Dynamics 365 e Microsoft 365, fundamentais para o avanço na carreira e reconhecimento profissional.
\end{itemize}
 % Completo e Validado
\section{Esta empresa possui programas sociais para as comunidades nas quais atua? Quais?}
A Petrobras demonstra um compromisso forte com o desenvolvimento social ao investir R\$ 200 milhões ao longo de quatro anos para treinar profissionalmente comunidades vulneráveis. Esse esforço é complementado por um orçamento significativo destinado a treinamentos internos, evidenciando o empenho da Petrobras no desenvolvimento contínuo de seus empregados e na promoção de avanços profissionais sustentáveis.
Além disso, a Petrobras desenvolve diversos programas sociais focados na educação, segurança, sustentabilidade ambiental e inclusão social, impactando as comunidades em suas áreas de atuação. Entre os principais programas estão:

\begin{itemize}
    \item \textbf{Programa Petrobras Socioambiental}: Este programa engloba projetos voltados para a educação, desenvolvimento econômico sustentável, e conservação ambiental, com a seleção de novos projetos nas áreas de cultura, esporte, negócios, ciência e tecnologia.
    \item \textbf{Iniciativa Territórios da Primeira Infância}: Foca na promoção dos direitos de crianças de 0 a 6 anos em 15 municípios nos estados de SP, RJ, ES, RN e SE, integrando projetos voltados para esta faixa etária.
    \item \textbf{Projeto Comunidade Segura}: Visa mitigar riscos para comunidades situadas nas faixas de dutos, fortalecendo organizações locais por meio de capacitação e consultoria, e promovendo o esporte educacional.
    \item \textbf{Tecendo as Águas}: Promove a preservação e uso sustentável dos recursos hídricos no litoral norte de São Paulo, integrando práticas de turismo responsável e educação ambiental.
    \item \textbf{Tecendo Redes na Primeira Infância}: Oferece suporte psicossocial e pedagógico a crianças na primeira infância, com acompanhamento interdisciplinar para prevenir e minimizar situações de violência.
\end{itemize}

 % Completo e Validado
\section{Esta empresa possui politica de preservação do meio ambiente e sustentabilidade? Quais?}

A Microsoft demonstra um forte compromisso com a sustentabilidade ambiental, adotando várias políticas e iniciativas:

\begin{itemize}
    \item \textbf{Carbono Negativo (Carbon Negative)}: A Microsoft se comprometeu a remover mais carbono da atmosfera do que emite até 2030. Este objetivo é alcançado através de uma combinação de redução de emissões, projetos de sequestro de carbono e investimentos em energia renovável.
    \item \textbf{Positiva em Água (Water Positive)}: A política visa a repor mais água do que a empresa consome em suas operações globais até 2030. Isso é realizado através de projetos de conservação e reabastecimento de água, além do uso eficiente da água nos processos operacionais.
    \item \textbf{Resíduos Zero (Zero Waste)}: A Microsoft busca eliminar os resíduos de suas operações diretas, produtos e embalagens até 2030. Isso inclui maximizar a reciclagem, apoiar a economia circular e inovar em soluções de embalagens sustentáveis.
    \item \textbf{Proteger e Preservar Ecossistemas (Protect and Preserve Ecosystems)}: A empresa está comprometida com a proteção e restauração de mais terra do que utiliza para suas operações, promovendo a biodiversidade e apoiando a saúde dos ecossistemas através de tecnologias e parcerias.
\end{itemize} % Completo e Validado

% ===============================================================
% Conclusão geral ===============================================
% \newpage
\setstretch{1.5} % Espaçamento entre linhas, certifique-se de que o pacote setspace está incluído em document.tex

% % Estilo de citação (plain, abbrv, ieeetr, etc.)
% \bibliographystyle{IEEEtran}  % Estilo que inclui URLs
% % \bibliographystyle{plain}

% \bibliography{
%     perguntas/1-foundation/latex/references,
%     perguntas/2-graduation/latex/references,
%     perguntas/3-main-productions/latex/references,
%     perguntas/4-location/latex/references,
%     perguntas/5-competitors/latex/references,
%     perguntas/6-protection/latex/references
% }

% Estilo de citação (plain, abbrv, ieeetr, etc.)
\bibliographystyle{IEEEtran}  % Estilo que inclui URLs
% \bibliographystyle{plain}

\bibliography{perguntas/1-foundation/latex/references,perguntas/2-graduation/latex/references,perguntas/3-main-products/latex/references.bib,perguntas/4-location/latex/references,perguntas/5-competitors/latex/references,perguntas/6-protection/latex/references, perguntas/7-profit/latex/references, perguntas/8-train/latex/references, perguntas/9-social-programs/latex/references, perguntas/10-green-purposes/latex/references}


% ===============================================================
% Referencias ===================================================
\newpage
\setstretch{1.5} % Espaçamento entre linhas, certifique-se de que o pacote setspace está incluído em document.tex

% % Estilo de citação (plain, abbrv, ieeetr, etc.)
% \bibliographystyle{IEEEtran}  % Estilo que inclui URLs
% % \bibliographystyle{plain}

% \bibliography{
%     perguntas/1-foundation/latex/references,
%     perguntas/2-graduation/latex/references,
%     perguntas/3-main-productions/latex/references,
%     perguntas/4-location/latex/references,
%     perguntas/5-competitors/latex/references,
%     perguntas/6-protection/latex/references
% }

% Estilo de citação (plain, abbrv, ieeetr, etc.)
\bibliographystyle{IEEEtran}  % Estilo que inclui URLs
% \bibliographystyle{plain}

\bibliography{perguntas/1-foundation/latex/references,perguntas/2-graduation/latex/references,perguntas/3-main-products/latex/references.bib,perguntas/4-location/latex/references,perguntas/5-competitors/latex/references,perguntas/6-protection/latex/references, perguntas/7-profit/latex/references, perguntas/8-train/latex/references, perguntas/9-social-programs/latex/references, perguntas/10-green-purposes/latex/references}


% ===============================================================
% Anexo ou Apendices ============================================
% \newpage
\setstretch{1.5} % Espaçamento entre linhas, certifique-se de que o pacote setspace está incluído em document.tex

% % Estilo de citação (plain, abbrv, ieeetr, etc.)
% \bibliographystyle{IEEEtran}  % Estilo que inclui URLs
% % \bibliographystyle{plain}

% \bibliography{
%     perguntas/1-foundation/latex/references,
%     perguntas/2-graduation/latex/references,
%     perguntas/3-main-productions/latex/references,
%     perguntas/4-location/latex/references,
%     perguntas/5-competitors/latex/references,
%     perguntas/6-protection/latex/references
% }

% Estilo de citação (plain, abbrv, ieeetr, etc.)
\bibliographystyle{IEEEtran}  % Estilo que inclui URLs
% \bibliographystyle{plain}

\bibliography{perguntas/1-foundation/latex/references,perguntas/2-graduation/latex/references,perguntas/3-main-products/latex/references.bib,perguntas/4-location/latex/references,perguntas/5-competitors/latex/references,perguntas/6-protection/latex/references, perguntas/7-profit/latex/references, perguntas/8-train/latex/references, perguntas/9-social-programs/latex/references, perguntas/10-green-purposes/latex/references}


% ===============================================================
% End Document ==================================================
\end{document}

